% =============================================================
%  ch04a_posenschaetzung.tex  --  6-DoF-Posenschätzung mit gelernten
%  Surface Embeddings (theoretische Grundlagen für Projekt 4)
% =============================================================

\chapter{6-DoF-Posenschätzung mit gelernten Surface Embeddings}

Dieses Kapitel vertieft die theoretischen Grundlagen hinter Projekt~4. Es baut
direkt auf dem Stereosehen (Kapitel zuvor) auf: Dort wurde die 6D-Pose über
klassische Merkmale (\code{PnP}, \code{ICP}) bestimmt; hier wird der moderne,
lernbasierte Weg über \textbf{dichte Surface Embeddings} (in der Linie von
\textbf{SurfEmb} \cite{haugaard2022surfemb}) entwickelt, der bei
texturarmen, glänzenden oder symmetrischen Industrieteilen robuster ist.

\section{Das Problem: sechs Freiheitsgrade}
\textbf{6-DoF-Posenschätzung} bestimmt \emph{Position und Orientierung} eines
bekannten Bauteils -- drei Translations- und drei Rotationsfreiheitsgrade.
\textbf{Warum überhaupt?} Ein Werkstück liegt selten zweimal gleich (Schüttgut
in Kisten, verschobene Vorrichtung, Förderband). Der Roboter muss das Teil
\emph{sehen}, bevor er greift. Gesucht ist die starre Transformation, die das
bekannte CAD-Modell in das Kamera-/Roboterkoordinatensystem überführt:
\begin{equation}
\boxed{\;\bm{p}_{\text{cam}} = \bm{R}\,\bm{p}_{\text{model}} + \bm{t}\;}
\qquad \bm{R}\in SO(3),\;\; \bm{t}\in\mathbb{R}^3
\end{equation}
Die rücktransformierte (``ausgerichtete'') Punktwolke ist nur die visuelle
Bestätigung -- das eigentliche Ergebnis ist $(\bm{R},\bm{t})$.

\section{Rotationsrepräsentationen}
Die Translation ist trivial ($\mathbb{R}^3$); die Rotation ist der heikle Teil,
weil $SO(3)$ gekrümmt ist. Gängige Darstellungen mit ihren Fallstricken:
\begin{longtable}{@{}p{3.6cm}p{10.6cm}@{}}
\toprule
\textbf{Darstellung} & \textbf{Eigenschaft / Fallstrick} \\
\midrule
Rotationsmatrix ($3{\times}3$) & 9 Zahlen, 6 Nebenbedingungen (orthonormal); redundant aber stabil \\
Euler-Winkel              & nur 3 Zahlen, aber \textbf{Gimbal Lock} (Freiheitsgrad-Verlust) \\
Quaternion                & 4 Zahlen, singularitätsfrei, aber \textbf{Double Cover} ($\bm{q}$ und $-\bm{q}$ = gleiche Drehung) \\
Kontinuierliche 6D-Form   & beste Wahl für die \textbf{NN-Regression} -- stetig, keine Sprünge \cite{zhou2019continuity} \\
\bottomrule
\end{longtable}
\noindent Für neuronale Netze ist die Stetigkeit entscheidend: Euler und
Quaternion haben Unstetigkeiten, an denen das Netz schlecht lernt. Die
6D-Darstellung (zwei Spaltenvektoren $\rightarrow$ Gram-Schmidt $\rightarrow$
$\bm{R}$) ist überall stetig.

\section{Warum gelernte Embeddings statt Regression}
Zwei naive Wege scheitern in der Industrie:
\begin{itemize}
\item \textbf{Handgefertigte Deskriptoren} (\code{FPFH} \cite{rusu2009fpfh},
  \code{SIFT}) versagen bei texturlosen, glänzenden, symmetrischen Teilen und
  unter Verdeckung -- genau die typischen Fälle in der Fertigung.
\item \textbf{Direkte Pose-Regression} (Netz $\rightarrow$ ein $(\bm{R},\bm{t})$)
  nimmt \emph{eine} Antwort an (unimodal) und \textbf{bricht bei Symmetrie}: ein
  symmetrisches Teil hat mehrere korrekte Posen, der Mittelwert ist falsch.
\end{itemize}
\textbf{Gelernte dichte Surface Embeddings} lösen beides: Ein Netz bildet jeden
Oberflächenpunkt in einen Merkmalsraum ab; gemessene Punkte werden über die
\emph{Nähe im Merkmalsraum} den Punkten des bekannten CAD-Modells zugeordnet.
Statt einer Punktschätzung lernt das Modell eine \textbf{Verteilung} über
Korrespondenzen -- Symmetrie wird implizit behandelt (mehrere Oberflächenpunkte
teilen sich dann schlicht ein Embedding) \cite{haugaard2022surfemb}.

\section{Kontrastives Lernen der Embeddings}
Trainiert werden die Embeddings mit einem \textbf{kontrastiven Verlust}: wahre
Pixel-$\leftrightarrow$-Punkt-Paare werden im Merkmalsraum zusammengezogen, alle
anderen Paare auseinandergedrückt. Die \emph{Positiv}-Paare sind die unter
bekannter Render-Pose bekannten Korrespondenzen, die \emph{Negativ}-Paare andere
gesampelte Oberflächenpunkte. In InfoNCE-Form für ein Query-Embedding $\bm{q}$
mit positivem Schlüssel $\bm{k}^+$:
\begin{equation}
\mathcal{L} = -\log
\frac{\exp\!\big(\mathrm{sim}(\bm{q},\bm{k}^+)/\tau\big)}
     {\sum_{j}\exp\!\big(\mathrm{sim}(\bm{q},\bm{k}_j)/\tau\big)}
\end{equation}
mit Kosinus-Ähnlichkeit $\mathrm{sim}(\cdot)$ und \textbf{Temperatur} $\tau$
(Tuning-Größe). \textbf{Der Clou:} die Labels sind quasi gratis -- man rendert
das CAD-Modell in bekannten Posen, womit Ground-Truth-Pose \emph{und}
Korrespondenzen automatisch vorliegen; keine manuellen Symmetrie-Labels nötig.

\section{Das Two-Tower-Modell (Query--Key)}
Die Architektur besteht aus zwei getrennten Türmen, die erst im gemeinsamen
Embedding-Raum zusammentreffen:
\begin{itemize}
\item \textbf{Query-Tower} -- \emph{Eingang:} Bildausschnitt (2D-Pixel, RGB/RGB-D);
  \emph{Ausgang:} Embedding pro Pixel; \emph{Architektur:} großes
  CNN-Encoder-Decoder-Netz (U-Net / ResNet).
\item \textbf{Key-Tower} -- \emph{Eingang:} eine 3D-Oberflächenkoordinate
  $(x,y,z)$ vom CAD; \emph{Ausgang:} Embedding dieses Punktes im selben Raum;
  \emph{Architektur:} kleines MLP (wenige Schichten reichen, da der Eingang
  simpel ist -- die \textbf{Asymmetrie} der Türme ist gewollt).
\end{itemize}
Beide werden \emph{gemeinsam} mit dem kontrastiven Verlust trainiert, sodass
zusammengehörige Pixel- und Punkt-Embeddings nahe beieinander landen. Ein Turm
spricht ``2D-Bild'', der andere ``3D-Oberfläche''; der gemeinsame Raum ist die
gemeinsame Sprache.

\begin{center}
\begin{tikzpicture}[node distance=6mm,every node/.style={font=\scriptsize}]
\node[node] (img) {Bildausschnitt\\2D-Pixel (RGB/RGB-D)};
\node[node,below=of img] (cnn) {CNN Encoder-Decoder\\U-Net / ResNet (groß)};
\node[node,below=of cnn] (pe) {Embedding pro Pixel};
\node[node,fill=accent2!12,draw=accent2,right=42mm of img] (pt) {3D-Oberflächenpunkt\\$(x,y,z)$ vom CAD};
\node[node,fill=accent2!12,draw=accent2,below=of pt] (mlp) {kleines MLP\\wenige Schichten};
\node[node,fill=accent2!12,draw=accent2,below=of mlp] (ke) {Embedding pro Punkt};
\node[node,fill=accent!18,below=12mm of $(pe)!0.5!(ke)$] (sh)
  {Gemeinsamer Embedding-Raum\\\itshape gemeinsam via kontrastivem Verlust trainiert};
\node[node,below=8mm of sh] (out)
  {Nächsten-Nachbarn-Matching\\dichte 2D-3D-Korrespondenzen $\rightarrow$ \code{PnP}+\code{RANSAC} $\rightarrow$ $(\bm{R},\bm{t})$};
\draw[flow] (img)--(cnn); \draw[flow] (cnn)--(pe);
\draw[flow] (pt)--(mlp);  \draw[flow] (mlp)--(ke);
\draw[flow] (pe)--(sh);   \draw[flow] (ke)--(sh);
\draw[flow] (sh)--(out);
\node[font=\scriptsize\bfseries,accent] at ($(cnn)+(0,1.0)$) {Query-Tower};
\node[font=\scriptsize\bfseries,accent2] at ($(mlp)+(0,1.0)$) {Key-Tower};
\end{tikzpicture}
\captionof{figure}{Two-Tower-Modell: die Türme bleiben getrennt bis zum
gemeinsamen Embedding-Raum -- genau diese Trennung \emph{ist} die Architektur.
Bei der Inferenz wird jedes Pixel-Embedding seinem nächsten
Oberflächenpunkt-Embedding zugeordnet.}
\end{center}

\section{Von Korrespondenzen zur Pose: PnP + RANSAC}
Aus den dichten 2D-3D-Korrespondenzen wird die Pose über \textbf{Perspective-n-Point}
(\code{PnP}) gelöst -- die $(\bm{R},\bm{t})$, die die 3D-Modellpunkte $\bm{P}_i$
auf die beobachteten Pixel $\bm{u}_i$ projiziert:
\begin{equation}
(\bm{R},\bm{t}) = \arg\min_{\bm{R},\bm{t}}
\sum_i \big\lVert \bm{u}_i - \pi\!\big(\bm{K}(\bm{R}\bm{P}_i + \bm{t})\big)\big\rVert^2
\end{equation}
mit Projektion $\pi(\cdot)$ und Intrinsik $\bm{K}$. \textbf{RANSAC} umhüllt das
Ganze und verwirft Ausreißer-Korrespondenzen (zieht zufällige minimale Mengen,
zählt Inlier, behält das beste Modell). Bei reinen 3D-3D-Daten tritt an die
Stelle der Reprojektion die euklidische Punktdistanz.

\section{Registrierung: grob und fein}
Eine \emph{einstufige} Registrierung ist unzuverlässig; Standard sind
\textbf{zwei Stufen}:
\begin{enumerate}
\item \textbf{Grob (global):} liefert die ``richtige Größenordnung'' aus
  Merkmalen -- Deskriptoren auf Quelle und Modell matchen, grobes
  $(\bm{R},\bm{t})$ lösen (z.\,B.\ \code{SAC-IA} + \code{FPFH}). Braucht
  \emph{keine} gute Startschätzung, ist aber nur näherungsweise.
\item \textbf{Fein = ICP} (\emph{Iterative Closest Point}
  \cite{besl1992icp}): zu jedem Quellpunkt den nächsten Zielpunkt suchen, die
  Gesamtdistanz minimieren, wiederholen bis zur Konvergenz:
  \begin{equation}
  E(\bm{R},\bm{t}) = \sum_i \big\lVert \bm{R}\,\bm{p}_i + \bm{t} - \bm{q}_{c(i)}\big\rVert^2,
  \qquad c(i)=\text{nächster Nachbar}
  \end{equation}
  Genau, aber nur \textbf{lokal}.
\end{enumerate}
\textbf{ICP-Schwächen:} (1) \emph{lokale Minima} -- es läuft bergab und bleibt
stecken; (2) \emph{initialisierungsabhängig} -- es funktioniert nur bei nahem
Start (daher läuft die Grobstufe zuerst). \textbf{Go-ICP}
\cite{yang2016goicp} garantiert das globale Optimum, indem es ganz $SE(3)$ (alle
starren Transformationen) per \emph{Branch-and-Bound} durchsucht: Raum
aufteilen, Fehler je Region beschränken, Regionen verwerfen, die nicht gewinnen
können -- langsamer, aber global optimal.

\section{Punktwolken-Netze: PointNet und PointNet++}
Verarbeitet man die Punktwolke direkt mit einem Netz, ist die zentrale
Schwierigkeit: eine Punktwolke ist eine \textbf{ungeordnete Menge} -- der Ausgang
darf nicht von der Reihenfolge abhängen (\textbf{Permutationsinvarianz}).
\begin{itemize}
\item \textbf{PointNet} \cite{qi2017pointnet}: dasselbe geteilte MLP auf jeden
  Punkt, dann \textbf{Max-Pooling} (symmetrisch = reihenfolgeunabhängig). Erfasst
  die \emph{globale} Form, verpasst aber lokale Details.
\item \textbf{PointNet++} \cite{qi2017pointnetpp}: wendet PointNet
  \emph{lokal} in Stufen an -- Zentren samplen (Farthest-Point-Sampling),
  Nachbarschaften gruppieren (Multi-Scale-Radien), lokale Merkmale extrahieren,
  auf größerer Skala wiederholen. Baut lokal + global hierarchisch auf, ähnlich
  einem CNN.
\end{itemize}

\section{Stereo, Tiefe und Skalierung}
\textbf{Warum Stereo?} Eine Mono-Kamera hat \emph{Skalenmehrdeutigkeit}
(klein-nah $\approx$ groß-fern). Stereo liefert Tiefe und löst die Skala -- über
Disparität, Basislinie und Brennweite (Triangulation, vgl.\ Stereo-Kapitel).
Zwei Wege in die Pipeline:
\begin{enumerate}[label=(\alph*)]
\item Tiefe als \textbf{Zusatzkanal} = RGB-D-Eingang, oder
\item Tiefe + Intrinsik \textbf{rückprojizieren} $\rightarrow$ Punktwolke
  $\rightarrow$ 3D-3D-Registrierung.
\end{enumerate}
\textbf{Vorverarbeitung} (vor dem Netz, ja): Kalibrierung (Intrinsik +
Extrinsik) und Rektifizierung, Tiefen-Entrauschung, Umgang mit ungültiger Tiefe
(glänzendes Metall liefert schlechte Tiefenwerte). \textbf{Auszahlung:} die
Tiefe gibt \emph{metrische} Skala -- $(\bm{R},\bm{t})$ in echten Millimetern, mit
denen der Roboter handeln kann.

\section{Zeitliche Konsistenz: Kalman-Pose-Tracking}
Eine bildweise geschätzte Pose ist verrauscht/zittrig. Ein \textbf{Kalman-Filter}
ergänzt \emph{zeitliche Konsistenz} (vgl.\ Kapitel zu Kalman/EKF): der
\emph{Zustand} ist die Pose (ggf.\ plus Geschwindigkeit), die \emph{Messung} ist
die bildweise Netz-Pose. Prädiktion der nächsten Pose $\rightarrow$ Korrektur mit
der neuen Schätzung. Das \textbf{glättet Jitter, verwirft Ausreißer-Frames und
überbrückt kurze Verdeckung / ausgefallene Detektionen} -- aus Einzelschuss wird
\emph{Tracking}, relevant für bewegtes Förderband oder armmontierte Kamera. Da
Rotationen nichtlinear sind, nutzt man EKF/UKF oder filtert auf einer geeigneten
Rotationsdarstellung.

\section{Sim-to-Real und Trainingsdaten}
Da die Labels durch Rendern des CAD-Modells praktisch gratis entstehen, wird
synthetisch trainiert -- aber es klafft ein \textbf{Sim-to-Real-Gap}
(synthetisch $\rightarrow$ real: Rauschen, Beleuchtung, Reflektanz). Gegenmittel:
\textbf{Domain Randomization}, physikalisch-basiertes Rendering (PBR, z.\,B.\ via
Isaac Sim, vgl.\ Kapitel zu Simulationswerkzeugen) und ein kleiner realer
Feintuning-Datensatz.

\section{Benchmarking und Metriken}
\begin{longtable}{@{}p{3.8cm}p{10.4cm}@{}}
\toprule
\textbf{Größe} & \textbf{Bedeutung} \\
\midrule
\textbf{ADD / ADD-S}     & mittlerer Punktabstand Modell$\leftrightarrow$Schätzung; ADD-S nutzt nächsten Punkt (für \emph{symmetrische} Teile) \\
Rotations-/Translationsfehler & geodätisch (Rotation) bzw.\ euklidisch (Translation), dazu RMSE \\
\textbf{BOP}-Metriken    & VSD, MSSD, MSPD $\rightarrow$ Average Recall \cite{hodan2018bop} \\
Industrie-KPIs           & Taktzeit/Latenz, Pick-Erfolgsrate, Falsch-Positiv-Rate \\
Datensätze               & BOP-Suite, speziell \textbf{T-LESS} und \textbf{ITODD} (texturlos/industriell) \\
\bottomrule
\end{longtable}

\begin{infobox}[title=Kernaussage in einem Satz]
6-DoF-Posenschätzung über gelernte Surface Embeddings ersetzt fragile
handgefertigte Deskriptoren durch ein \emph{Zwei-Turm-Netz}, das Pixel und
CAD-Oberflächenpunkte in einen gemeinsamen Merkmalsraum abbildet; das dichte
Korrespondenz-Matching liefert über \code{PnP}+\code{RANSAC} (optional
ICP-verfeinert) die starre Transformation $(\bm{R},\bm{t})$ -- robust auch bei
texturlosen, glänzenden und symmetrischen Teilen.
\end{infobox}
